\documentclass[tikz,border=6pt]{standalone}
\usepackage{peabra}
\begin{document}
\begin{tikzpicture}[x=1cm,y=1cm]
  \def\r{1.25}\def\w{0.66}
  \foreach \i/\szin/\txt in {
      8/peSargaH/{lakossági vészhelyzeti reagálás},
      7/peSargaH/{üzemi vészhelyzeti reagálás},
      6/peKiemelH/{kibocsátás utáni fizikai védelem (gát, felfogó medence)},
      5/peKiemelH/{fizikai védelem (hasadótárcsa, biztonsági szelep)},
      4/peZoldH/{biztonsági műszeres rendszer (SIS)},
      3/peZoldH/{kritikus riasztások és kezelői beavatkozás},
      2/peHalvany/{alapfolyamat-irányítás (BPCS)}} {
    \pgfmathsetmacro{\R}{\r+(\i-1)*\w}
    \fill[\szin, draw=white, line width=1.2pt] (-\R,0) arc (180:0:\R) -- cycle;
  }
  \fill[peKek, draw=white, line width=1.2pt] (-\r,0) arc (180:0:\r) -- cycle;
  \node[font=\scriptsize\bfseries, text=white, align=center] at (0,0.5) {folyamat-\\tervezés};
  \foreach \i/\txt in {2/{alapfolyamat-irányítás (BPCS)}, 3/{riasztások és kezelői beavatkozás}, 4/{biztonsági műszeres rendszer (SIS)},
      5/{fizikai védelem (hasadótárcsa, szelep)}, 6/{kibocsátás utáni védelem (gát, felfogó)}, 7/{üzemi vészhelyzeti reagálás}, 8/{lakossági vészhelyzeti reagálás}} {
    \pgfmathsetmacro{\R}{\r+(\i-1.5)*\w}
    \node[font=\scriptsize, text=peSzoveg] at (0,\R) {\txt};
  }
  \draw[peSotet, line width=0.8pt] (-6.2,0) -- (6.2,0);
  % csoportok jobbra
  \def\br#1#2#3#4{\pgfmathsetmacro{\a}{\r+(#1-2)*\w+0.04}\pgfmathsetmacro{\b}{\r+(#2-1)*\w-0.04}
     \draw[decorate, decoration={brace, amplitude=4pt}, #4, line width=0.8pt] (6.0,\b) -- (6.0,\a); \node[font=\scriptsize\bfseries, text=#4, anchor=west, align=left] at (6.2,{(\a+\b)/2}) {#3};}
  \br{2}{4}{megelőzés}{peZold}
  \br{5}{6}{következmény-\\csökkentés}{peKiemel}
  \br{7}{8}{vészhelyzeti\\reagálás}{peSarga}
  \node[font=\tiny, text=peSzurke, anchor=north] at (0,-0.1) {a LOPA a független védelmi rétegeket (IPL) számítja be a kockázatcsökkentésbe};
\end{tikzpicture}
\end{document}
